% =====================================================================
%  RBE 502 - Quadrotor Control Project: Geometric Tracking Control
%  Report (IEEEtran conference).
%  Compile from this folder: pdflatex skate_quadrotor_project.tex (x2)
%  Figures: figures/ (copied from the run folders and the plotting scripts).
% =====================================================================
\documentclass[conference]{IEEEtran}
\usepackage[utf8]{inputenc}
\usepackage{cite}
\usepackage{amsmath,amssymb,amsfonts}
\usepackage{graphicx}
\usepackage{xcolor}
\usepackage{booktabs}
\usepackage{multirow}
\usepackage{subcaption}
\captionsetup[subfigure]{labelformat=simple}
\renewcommand{\thesubfigure}{(\alph{subfigure})}
\usepackage{url}
\usepackage{stfloats}   % fixes ordering of full-width figures

\setcounter{dbltopnumber}{3}
\renewcommand{\dbltopfraction}{0.95}
\renewcommand{\dblfloatpagefraction}{0.65}
\renewcommand{\topfraction}{0.9}
\renewcommand{\floatpagefraction}{0.7}
\renewcommand{\textfraction}{0.05}
\setlength{\dblfloatsep}{8pt plus 2pt minus 2pt}
\setlength{\dbltextfloatsep}{10pt plus 2pt minus 2pt}
\setlength{\textfloatsep}{8pt plus 2pt minus 2pt}

% \fig{file}{height}: includes figures/file if present, else a placeholder
\newcommand{\fig}[2]{%
  \IfFileExists{figures/#1}%
    {\includegraphics[width=\linewidth]{figures/#1}}%
    {\fbox{\parbox[c][#2][c]{0.92\linewidth}{\centering\scriptsize\texttt{figures/\detokenize{#1}}}}}}

% ---------- Notation -------------------------------------------------
\newcommand{\vx}{\mathbf{x}}
\newcommand{\vv}{\mathbf{v}}
\newcommand{\va}{\mathbf{a}}
\newcommand{\ve}{\mathbf{e}}
\newcommand{\vb}{\mathbf{b}}
\newcommand{\R}{\mathbf{R}}
\newcommand{\J}{\mathbf{J}}
\newcommand{\M}{\mathbf{M}}
\newcommand{\F}{\mathbf{F}}
\newcommand{\K}{\mathbf{K}}
\newcommand{\Om}{\boldsymbol{\Omega}}

\begin{document}

\title{Geometric Tracking Control of a Crazyflie 2.X Quadrotor on SE(3)}

\author{
\IEEEauthorblockN{Shreehan Santosh Kate}
\IEEEauthorblockA{Robotics Engineering\\Worcester Polytechnic Institute\\skate@wpi.edu}
}

\maketitle

% =====================================================================
\begin{abstract}
We implement the geometric tracking controller of Lee et al.\ for a simulated
Crazyflie 2.X and use it to track a circle and a diamond trajectory. The
gains are designed from second-order models of the attitude and position
loops and then tuned with a sweep of the position-loop natural frequency.
With the final gains, the position error RMSE is 2.6~mm on the circle and
4.2~mm on the diamond. Halving the derivative gain makes the diamond
oscillate at the loop's natural frequency and triples its error, while
halving the proportional gain mainly increases the error on the circle.
Shortening the trajectories increases the error much faster than the speed.
The two trajectories reach different limits: on the diamond, the attitude
loop cannot follow the required tilt rate, and on the circle, the rotors
reach their maximum thrust.
\end{abstract}

\begin{IEEEkeywords}
Quadrotor, geometric control, SE(3), differential flatness, gain tuning
\end{IEEEkeywords}

% =====================================================================
\section{Introduction}
Controllers obtained by linearizing the quadrotor dynamics about hover assume
small roll and pitch angles. The geometric controller of~\cite{lee} instead
defines the attitude error directly on the rotation matrix, so it remains
valid far from hover and has no Euler-angle singularities. This report
presents its implementation for the Crazyflie 2.X in a PyBullet simulator,
the design of its gains (Section~\ref{sec:best}), the sensitivity of the
tracking to the position gains (Section~\ref{sec:sens}), and the effect of
the trajectory speed (Section~\ref{sec:speed}).

% =====================================================================
\section{Model and Controller}\label{sec:model}
The inertial frame has $\ve_3$ pointing up. The state is the position $\vx$,
the velocity $\vv$, the rotation matrix $\R=[\vb_1\ \vb_2\ \vb_3]$, and the
body angular velocity $\Om$. With mass $m=0.027$~kg and inertia
$\J=\mathrm{diag}(1.4,1.4,2.17)\times10^{-5}$~kg\,m$^2$,
\begin{equation}
m\dot\vv = f\R\ve_3 - mg\ve_3,\qquad \J\dot\Om + \Om\times\J\Om = \M.
\label{eq:dyn}
\end{equation}
The thrust $f$ acts only along $\vb_3$, so the quadrotor must tilt to
accelerate horizontally. This leads to a cascaded controller~\cite{mahony}.

\emph{Position loop.} With $\ve_x=\vx-\vx_d$ and $\ve_v=\vv-\vv_d$,
\begin{equation}
\F_d = -\K_x\ve_x - \K_v\ve_v + mg\ve_3 + m\va_d,\qquad f = \F_d\cdot\R\ve_3 .
\label{eq:pos}
\end{equation}
The desired attitude aligns $\vb_3$ with $\F_d$ and sets the heading from the
desired yaw $\psi_d$: $\vb_{3,d}=\F_d/\|\F_d\|$,
$\vb_{1,c}=[\cos\psi_d\ \sin\psi_d\ 0]^\top$,
$\vb_{2,d}=\vb_{3,d}\times\vb_{1,c}/\|\vb_{3,d}\times\vb_{1,c}\|$,
$\vb_{1,d}=\vb_{2,d}\times\vb_{3,d}$, and
$\R_d=[\vb_{1,d}\ \vb_{2,d}\ \vb_{3,d}]$.

\emph{Attitude loop.} With $\Om_d=\mathbf 0$,
\begin{equation}
\ve_R=\tfrac12\left(\R_d^\top\R-\R^\top\R_d\right)^\vee,\quad
\M=-\K_R\ve_R-\K_\Omega\Om+\Om\times\J\Om .
\label{eq:att}
\end{equation}
The simulator integrates the dynamics at 1000~Hz, and the controller runs at
200~Hz. The provided code converts $f$ and $\M$ into rotor speeds and clips
each rotor thrust to $[0,\ 0.149]$~N, so the maximum total thrust is
$f_{\max}=2.25\,mg$.

% =====================================================================
\begin{figure*}[!t]
\centering
\begin{subfigure}{0.32\textwidth}\centering
  \fig{task1_circle_position.png}{4cm}\caption{Position vs.\ desired}\end{subfigure}\hfill
\begin{subfigure}{0.32\textwidth}\centering
  \fig{task1_circle_error.png}{4cm}\caption{Position error $\ve_x$}\end{subfigure}\hfill
\begin{subfigure}{0.32\textwidth}\centering
  \fig{task1_circle_path3d.png}{4cm}\caption{3D path vs.\ desired}\end{subfigure}\\[4pt]
\begin{subfigure}{0.24\textwidth}\centering
  \fig{task1_circle_attitude.png}{5cm}\caption{Roll, pitch, yaw of $\R$, $\R_d$; $\|\ve_R\|$}\end{subfigure}\hspace{0.05\textwidth}
\begin{subfigure}{0.24\textwidth}\centering
  \fig{task1_circle_control.png}{5cm}\caption{Thrust $f$ and moment $\M$}\end{subfigure}
\caption{Circle (15~s) with the final gains. In (c), the actual path covers
the desired path.}
\label{fig:t1_circle}
\end{figure*}

\section{Best Gains Performance}\label{sec:best}

\subsection{Gain Design}
For small errors, $\ve_R$ approximates the roll, pitch, and yaw errors, and
the term $\Om\times\J\Om$ in~\eqref{eq:att} cancels the gyroscopic coupling
in~\eqref{eq:dyn}. Each attitude axis then behaves as
$J_i\ddot\theta+k_{\Omega,i}\dot\theta+k_{R,i}\theta=0$. Assuming that the
attitude loop is fast, each position axis behaves as
$m\ddot e+k_v\dot e+k_xe=0$. The gains therefore follow from a chosen
natural frequency $\omega_n$ and damping ratio $\zeta$:
\begin{equation}
k_R=J_i\omega_n^2,\ k_\Omega=2\zeta J_i\omega_n,\qquad
k_x=m\omega_n^2,\ k_v=2\zeta m\omega_n .
\end{equation}
We use $\zeta=0.9$ for both loops. For roll and pitch, $\omega_n=25$~rad/s
gives a settling time $4/(\zeta\omega_n)\approx0.18$~s and
$\omega_nT_s=0.125$ at the 5~ms control period. The yaw moment comes only
from the rotor reaction torques and is limited to about
$\gamma mg\approx6.6\times10^{-3}$~N\,m. With $\omega_n=25$~rad/s, the
initial yaw error of the hover test would already require this moment, so we
use $\omega_n=15$~rad/s for yaw. Yaw does not affect $\vb_3$, so the slower
yaw loop does not slow down position tracking. Table~\ref{tab:gains} lists
the final gains.

\begin{table}[htbp]
\caption{Final Controller Gains}
\label{tab:gains}
\centering
\setlength{\tabcolsep}{4pt}
\begin{tabular}{lll}
\toprule
Gain & Value & Design ($\zeta=0.9$)\\
\midrule
$\K_x$ & $2.700\,\mathbf I_3$ & $\omega_n=10$ rad/s\\
$\K_v$ & $0.486\,\mathbf I_3$ & $\omega_n=10$ rad/s\\
$\K_R$ & $\mathrm{diag}(8.75,\,8.75,\,4.88)\times10^{-3}$ & $\omega_n=25,25,15$ rad/s\\
$\K_\Omega$ & $\mathrm{diag}(6.30,\,6.30,\,5.86)\times10^{-4}$ & $\omega_n=25,25,15$ rad/s\\
\bottomrule
\end{tabular}
\end{table}

\emph{Hover test.} Starting from roll $20^\circ$, pitch $-15^\circ$, and yaw
$30^\circ$ (attitude error $36.8^\circ$), the attitude error falls below
$1^\circ$ after 0.29~s. This matches the predicted yaw settling time of
0.30~s, since yaw is the slowest axis. From a $45^\circ$ roll it converges in
0.20~s (predicted 0.18~s). The smallest rotor thrust is 0.0036~N and
0.011~N in these two tests, so the rotors never reach their lower limit.

\begin{figure*}[!t]
\centering
\begin{subfigure}{0.32\textwidth}\centering
  \fig{task1_diamond_position.png}{4cm}\caption{Position vs.\ desired}\end{subfigure}\hfill
\begin{subfigure}{0.32\textwidth}\centering
  \fig{task1_diamond_error.png}{4cm}\caption{Position error $\ve_x$}\end{subfigure}\hfill
\begin{subfigure}{0.32\textwidth}\centering
  \fig{task1_diamond_path3d.png}{4cm}\caption{3D path vs.\ desired}\end{subfigure}\\[4pt]
\begin{subfigure}{0.24\textwidth}\centering
  \fig{task1_diamond_attitude.png}{5cm}\caption{Roll, pitch, yaw of $\R$, $\R_d$; $\|\ve_R\|$}\end{subfigure}\hspace{0.05\textwidth}
\begin{subfigure}{0.24\textwidth}\centering
  \fig{task1_diamond_control.png}{5cm}\caption{Thrust $f$ and moment $\M$}\end{subfigure}
\caption{Diamond (8~s) with the final gains. The $e_z$ peaks at the start and
end come from the 12.5~mm ground offset.}
\label{fig:t1_diamond}
\end{figure*}

\subsection{Tuning of the Position Loop}
With the attitude gains fixed, we varied the position-loop $\omega_n$ with
$\zeta=0.9$ (Table~\ref{tab:sweep}). On the circle, the error decreases
approximately as $\omega_n^{-1.8}$ up to 12~rad/s. On the diamond, it
reaches its minimum at $\omega_n=8$--10~rad/s and increases at 12~rad/s,
where the peak attitude error reaches $17.7^\circ$ and the circle error
starts to oscillate. At 14~rad/s, both trajectories become unstable and the
quadrotor flips over. As the position loop approaches the 25~rad/s attitude
loop, the attitude can no longer follow $\R_d$, which is the time-scale
separation that the cascade assumes. No rotor saturates at 10 or 12~rad/s,
so saturation is not the cause. We chose $\omega_n=10$~rad/s: it gives the
lowest diamond error and nearly the lowest combined error, and it keeps a
ratio of 2.5 to the attitude loop instead of 2.1 at 12~rad/s.

\begin{table}[htbp]
\caption{Position-Loop Sweep ($\zeta=0.9$). RMSE$_e$ in mm; $\theta_{\max}$: peak
attitude error angle on the diamond}
\label{tab:sweep}
\centering
\setlength{\tabcolsep}{5pt}
\begin{tabular}{cccccc}
\toprule
$\omega_n$ (rad/s) & $k_x$ & $k_v$ & Circle & Diamond & $\theta_{\max}$\\
\midrule
4 & 0.432 & 0.194 & 12.81 & 7.39 & $4.1^\circ$\\
6 & 0.972 & 0.292 & 6.62 & 4.73 & $3.9^\circ$\\
8 & 1.728 & 0.389 & 3.94 & 4.20 & $7.6^\circ$\\
\textbf{10} & \textbf{2.700} & \textbf{0.486} & \textbf{2.59} & \textbf{4.17} & $12.5^\circ$\\
12 & 3.888 & 0.583 & 1.85 & 4.99 & $17.7^\circ$\\
14 & 5.292 & 0.680 & \multicolumn{3}{c}{unstable (flips)}\\
\bottomrule
\end{tabular}
\end{table}

\begin{figure*}[!t]
\centering
\begin{subfigure}{0.42\textwidth}\centering
  \fig{task2_circle.pdf}{6cm}\caption{Circle}\end{subfigure}\hfill
\begin{subfigure}{0.42\textwidth}\centering
  \fig{task2_diamond.pdf}{6cm}\caption{Diamond}\end{subfigure}
\caption{Position error components for the baseline gains and with $\K_v$ or
$\K_x$ reduced by 50\%.}
\label{fig:sens}
\end{figure*}

\subsection{Tracking Results}
Table~\ref{tab:task1} lists the RMSE with the final gains, and
Figs.~\ref{fig:t1_circle} and~\ref{fig:t1_diamond} show the responses. Both
trajectories score below the 0.05~m threshold by about an order of
magnitude. The acceleration feedforward $m\va_d$ supplies almost all of the
required thrust, and the actual thrust stays close to the ideal thrust of
Section~\ref{sec:speed} (Fig.~\ref{fig:t3_thrust}). On the circle, the
horizontal errors are about six times the vertical error, because
horizontal acceleration requires the quadrotor to tilt first, whereas
vertical acceleration only changes the thrust. The circle requires a tilt of
about $30^\circ$ at its highest speed (2.36~m/s, centripetal acceleration
$5.5$~m/s$^2$). The diamond starts and ends at $z=0$, but the center of mass
rests 12.5~mm above the ground, so the start and end of every run contain an
error that the controller cannot remove. Without the first and last 0.5~s,
the diamond RMSE$_e$ is 1.76~mm.

\begin{table}[htbp]
\caption{Tracking Performance with the Final Gains (mm)}
\label{tab:task1}
\centering
\setlength{\tabcolsep}{4pt}
\begin{tabular}{lccccc}
\toprule
& \multicolumn{4}{c}{RMSE} & Max\\
\cmidrule(lr){2-5}
Trajectory & $x$ & $y$ & $z$ & $e$ & $\|\ve_x\|$\\
\midrule
Circle (15 s) & 1.78 & 1.87 & 0.30 & 2.59 & 7.2\\
Diamond (8 s) & 0.10 & 2.81 & 3.07 & 4.17 & 16.5\\
\bottomrule
\end{tabular}
\end{table}

% =====================================================================
\begin{figure*}[!t]
\centering
\begin{subfigure}{0.48\textwidth}\centering
  \fig{task3_circle_fd.pdf}{4.5cm}\caption{Circle}\end{subfigure}\hfill
\begin{subfigure}{0.48\textwidth}\centering
  \fig{task3_diamond_fd.pdf}{4.5cm}\caption{Diamond}\end{subfigure}
\caption{Required thrust $f_d(t)$~\eqref{eq:fd} for the three durations,
with $f_{\max}=2.25\,mg$.}
\label{fig:fd}
\end{figure*}

\section{Sensitivity to Gain Variations}\label{sec:sens}
Starting from the final gains, we reduced $\K_v$ and then $\K_x$ by 50\%,
one at a time. Table~\ref{tab:sens} lists the results, and
Fig.~\ref{fig:sens} compares the errors with the baseline.

\begin{table}[htbp]
\caption{Gain Sensitivity: RMSE (mm)}
\label{tab:sens}
\centering
\setlength{\tabcolsep}{3.5pt}
\begin{tabular}{llcccccc}
\toprule
& Case & $\omega_n$ & $\zeta$ & $x$ & $y$ & $z$ & $e$\\
\midrule
\multirow{3}{*}{\rotatebox{90}{Circle}}
& Baseline & 10.0 & 0.90 & 1.78 & 1.87 & 0.30 & 2.59\\
& $\K_v\times0.5$ & 10.0 & 0.45 & 2.05 & 2.23 & 0.33 & 3.05\\
& $\K_x\times0.5$ & 7.07 & 1.27 & 2.79 & 2.96 & 0.55 & 4.11\\
\midrule
\multirow{3}{*}{\rotatebox{90}{Diamond}}
& Baseline & 10.0 & 0.90 & 0.10 & 2.81 & 3.07 & 4.17\\
& $\K_v\times0.5$ & 10.0 & 0.45 & 0.31 & 13.67 & 3.50 & 14.12\\
& $\K_x\times0.5$ & 7.07 & 1.27 & 0.17 & 2.35 & 3.47 & 4.19\\
\bottomrule
\end{tabular}
\end{table}

\emph{Derivative gain.} Halving $\K_v$ keeps $\omega_n$ and halves $\zeta$
to 0.45. The error does not grow through lag; it oscillates. On the diamond,
$e_y$ oscillates at 1.75~Hz (11~rad/s, close to $\omega_n$) with an
amplitude of 15--25~mm for the whole trajectory, and RMSE$_e$ increases
3.4 times. Each stop at a vertex excites this oscillation again. The
oscillation is larger than $\zeta=0.45$ alone suggests, because the lag of
the attitude loop reduces the effective damping of the position loop. On
the smooth circle, RMSE$_e$ increases by only 18\%, but the number of sign
changes of $e_x$ rises from 6 to 41, which shows a small persistent
oscillation.

\emph{Proportional gain.} Halving $\K_x$ lowers $\omega_n$ to 7.07~rad/s and
raises $\zeta$ to 1.27. The response becomes overdamped: the errors grow but
keep their shape, with no added oscillation. The circle error increases by
59\%, consistent with the trend of Table~\ref{tab:sweep}. The diamond error
is unchanged, and the peak attitude error decreases from $12.5^\circ$ to
$6.1^\circ$, because the softer position loop gives the attitude loop more
time.

The proportional gain therefore sets the tracking accuracy on sustained,
smooth motion, and the derivative gain sets the damping that abrupt,
stop-and-go motion requires. The RMSE alone can hide this difference, as the
circle with reduced $\K_v$ shows.

% =====================================================================
\begin{figure*}[!t]
\centering
\begin{subfigure}{0.42\textwidth}\centering
  \fig{task3_circle_errors.pdf}{6cm}\caption{Circle}\end{subfigure}\hfill
\begin{subfigure}{0.42\textwidth}\centering
  \fig{task3_diamond_errors.pdf}{6cm}\caption{Diamond}\end{subfigure}
\caption{Position error components for the three durations.}
\label{fig:t3_err}
\end{figure*}

\section{Effect of Trajectory Speed}\label{sec:speed}

\subsection{Required Thrust from Differential Flatness}
The position and yaw are flat outputs of the quadrotor~\cite{mellinger}. If the quadrotor tracks the trajectory exactly, then $\dot\vv=\va_d$, and
the translational dynamics~\eqref{eq:dyn} give
$f\R\ve_3=m(\va_d+g\ve_3)$. Since $\R\ve_3=\vb_3$ is a unit vector,
\begin{equation}
f_d(t)=m\,\|\va_d(t)+g\ve_3\|,\qquad
\vb_{3,d}=\frac{\va_d+g\ve_3}{\|\va_d+g\ve_3\|}.
\label{eq:fd}
\end{equation}
The desired acceleration alone thus determines the required thrust and the
tilt $\theta_d$, with $\cos\theta_d=(a_{d,z}+g)/\|\va_d+g\ve_3\|$. Both
trajectories consist of rest-to-rest quintic segments. Shortening the
duration from $T_0$ to $T$ therefore scales $\va_d$ by $s^2$, where
$s=T_0/T$, and $f_d=m\|s^2\va_{d,0}+g\ve_3\|$. At low speed, gravity
dominates and $f_d$ grows slowly. At high speed, the $s^2$ term dominates.

\begin{figure*}[!t]
\centering
\begin{subfigure}{0.42\textwidth}\centering
  \fig{task3_circle_thrust.pdf}{6cm}\caption{Circle}\end{subfigure}\hfill
\begin{subfigure}{0.42\textwidth}\centering
  \fig{task3_diamond_thrust.pdf}{6cm}\caption{Diamond}\end{subfigure}
\caption{Actual thrust $f(t)$ and required thrust $f_d(t)$ for the three
durations.}
\label{fig:t3_thrust}
\end{figure*}

\subsection{Choice of Durations}
For the shortest duration, the peak of $f_d$ must lie between $1.5\,mg$ and
$2.25\,mg$. We chose 8, 5, and 3.5~s for the diamond (peak $1.64\,mg$) and
15, 12, and 9~s for the circle (peak $1.86\,mg$), the shortest durations in
this range for which the quadrotor remained stable
(Table~\ref{tab:speed}). The provided circle has fixed segment times, so we
generated the shorter circles by time scaling,
$\vx_d(t)=\vx_0(st)$, $\vv_d=s\vv_0(st)$, $\va_d=s^2\va_0(st)$. This is
exact for quintic segments, and at $T=T_0$ it reproduces the default
circle. Fig.~\ref{fig:fd} shows $f_d(t)$.

\begin{figure*}[!t]
\centering
\begin{subfigure}{0.495\textwidth}\centering
  \fig{task3_circle_paths.pdf}{5cm}\caption{Circle}\end{subfigure}\hfill
\begin{subfigure}{0.495\textwidth}\centering
  \fig{task3_diamond_paths.pdf}{5cm}\caption{Diamond}\end{subfigure}
\caption{Actual and desired paths for the three durations: 3D view and
projection onto the plane of motion.}
\label{fig:t3_paths}
\end{figure*}

\subsection{Results}
Table~\ref{tab:speed} and Figs.~\ref{fig:t3_err}--\ref{fig:t3_paths} show the
results. In both cases the error grows much faster than the speed. The
diamond is 2.3 times faster at 3.5~s and its RMSE$_e$ is 6.2 times larger.
The circle is 1.7 times faster at 9~s and its RMSE$_e$ is 5.3 times larger.

\begin{table}[htbp]
\caption{Effect of Trajectory Duration}
\label{tab:speed}
\centering
\setlength{\tabcolsep}{2.6pt}
\begin{tabular}{lcccccccc}
\toprule
& $T$ & \multicolumn{2}{c}{Peak $f$ ($mg$)} & Peak & \multicolumn{4}{c}{RMSE (mm)}\\
\cmidrule(lr){3-4}\cmidrule(lr){6-9}
& (s) & $f_d$ & actual & tilt & $x$ & $y$ & $z$ & $e$\\
\midrule
\multirow{4}{*}{\rotatebox{90}{Circle}}
& 15 & 1.15 & 1.15 & $30^\circ$ & 1.78 & 1.87 & 0.30 & 2.59\\
& 12 & 1.34 & 1.35 & $42^\circ$ & 3.64 & 3.85 & 0.76 & 5.35\\
& 9 & 1.86 & 1.93 & $61^\circ$ & 9.26 & 9.72 & 2.50 & 13.65\\
& 8 & 2.23 & \multicolumn{6}{c}{unstable}\\
\midrule
\multirow{4}{*}{\rotatebox{90}{Diamond}}
& 8 & 1.11 & 1.14 & $14^\circ$ & 0.10 & 2.81 & 3.07 & 4.17\\
& 5 & 1.29 & 1.33 & $36^\circ$ & 0.22 & 6.75 & 3.73 & 7.72\\
& 3.5 & 1.64 & 1.88 & $86^\circ$ & 0.76 & 23.78 & 9.96 & 25.80\\
& 3.25 & 1.75 & \multicolumn{6}{c}{unstable}\\
\bottomrule
\end{tabular}
\end{table}

\subsection{Discussion}
\emph{Diamond: attitude rate.} At 3.5~s, the required tilt changes from
$0^\circ$ to $50^\circ$ and back on every edge within a few tenths of a
second. With $\Om_d=\mathbf 0$, the damping term $-\K_\Omega\Om$ opposes
this rotation as well, so the attitude lags $\R_d$ by up to $49^\circ$. The
resulting position error increases $\F_d$, the commanded tilt overshoots to
$86^\circ$, and the actual thrust rises to $1.88\,mg$ against
$f_d=1.64\,mg$ (Fig.~\ref{fig:t3_thrust}). When the quadrotor is tilted by
$\theta$, only $f\cos\theta$ acts vertically, so $e_z$ grows to 10~mm. The
smallest rotor thrust falls to 0.016~N, close to the lower limit, which
leaves little margin for the moment. The path cuts the vertices
(Fig.~\ref{fig:t3_paths}). At 3.25~s, the quadrotor flips although
$f_d=1.75\,mg$ remains below $f_{\max}$, so the total thrust is not the
limiting factor on this trajectory.

\emph{Circle: rotor thrust.} The circle's acceleration changes smoothly, so
the actual thrust follows $f_d$ closely even at 9~s ($1.93\,mg$ against
$1.86\,mg$). The required thrust peaks once, during the revolution. At 9~s,
the largest rotor thrust reaches 0.138~N of the 0.149~N limit. The
quadrotor lags behind the centripetal acceleration, so the path bulges
outward by up to 43~mm. At 8.5~s, the rotors saturate at their maximum
during the revolution, and at 8~s ($f_d=2.23\,mg$) the quadrotor becomes
unstable.

Faster trajectories therefore increase the error in three ways: they require
larger tilts, which reduce the vertical thrust; they require faster changes
of the tilt, which the attitude loop with $\Om_d=\mathbf 0$ follows with a
lag; and they leave less rotor thrust for the moment, both near zero thrust
(diamond) and near the maximum (circle). Feedforward of the desired angular
velocity, obtained from the jerk by differential flatness, would remove most
of the attitude lag.

% =====================================================================
\section{Conclusion}
The geometric controller, with gains designed from second-order models,
tracks both trajectories with an RMSE$_e$ of 2.6 and 4.2~mm. The
position-loop natural frequency is bounded by the attitude loop: the system
becomes unstable when their ratio falls below about two. Reducing the
derivative gain produces oscillation, and reducing the proportional gain
produces lag. When the trajectories are shortened, the diamond is limited by
the attitude rate and the circle by the rotor thrust, so the total thrust
alone does not determine whether a trajectory can be tracked.

% =====================================================================
\begin{thebibliography}{00}
\bibitem{lee} T. Lee, M. Leok, and N. H. McClamroch, ``Geometric tracking
control of a quadrotor UAV on SE(3),'' in \textit{Proc. IEEE Conf. Decision
and Control (CDC)}, 2010, pp. 5420--5425.
\bibitem{mahony} R. Mahony, V. Kumar, and P. Corke, ``Multirotor aerial
vehicles: Modeling, estimation, and control of quadrotor,'' \textit{IEEE
Robot. Autom. Mag.}, vol. 19, no. 3, pp. 20--32, 2012.
\bibitem{mellinger} D. Mellinger and V. Kumar, ``Minimum snap trajectory
generation and control for quadrotors,'' in \textit{Proc. IEEE Int. Conf.
Robot. Autom. (ICRA)}, 2011, pp. 2520--2525.
\end{thebibliography}
\end{document}
